\section{Related Work}
\label{sec:related}

\paragraph{Latent CoT and the Illusion of Superposition.}
COCONUT \citep{hao2024coconut} replaces explicit rationales with
continuous latent tokens; CODI \citep{shen2025codi} compresses CoT
into latents via self-distillation. \citet{rizvi2026illusion} report
that a fine-tuned COCONUT reaches $96.6\%$ on ProsQA \emph{without}
latents, $99.0\%$ with them, and $85.3\%$ with explicit CoT --- the
latent machinery is not doing measurable work in their fine-tuned
setting. Our vanilla SFT baseline at GPT-2 small reaches $81.77\%$,
roughly $15$pp below their $96.6\%$; we did not close this gap. The
qualitative phenomenon they report --- fine-tuned latent CoT models
reaching comparable accuracy without their latents --- replicates
at our operating point: matrix-CODI at $82.03\%$ vs.\ pure SFT at
$81.77\%$ (gap $0.26$pp, within three-seed noise). Our contribution
relative to that work is a structural argument about the training
objective that does not depend on matching their accuracy: the
matrix-bottleneck objective produces rank-indifferent gradients,
and four readouts nonlinear in $\Z$ are unable to bend the
rank-$k$ curve.

\paragraph{SIM-CoT.}
\citet{shen2025simcot} diagnose latent CoT instability as
insufficient step-level supervision and propose injecting per-step
targets. Our diagnosis is at a different layer (the matrix
bottleneck's objective produces rank-indifferent gradients,
adjudicated by four positive-control readouts). The two diagnoses
are consistent with both mechanisms operating simultaneously.

\paragraph{Reasoning by Superposition and CoT2.}
\citet{zhu2025superposition} prove that a two-layer transformer
with $D$ steps of continuous thought can solve directed graph
reachability, with each thought encoding a parallel BFS frontier.
\citet{gozeten2025cot2} show similar parallel-exploration behavior
under a GRPO-style training regime. Both are theoretical capacity
results with small empirical demonstrations. Capacity and
what-gets-learned are distinct; our result is about what CODI
distillation shapes, not whether transformers can in principle
encode superposition.

\paragraph{February 2026 rank measurements (direct adjacency).}
\citet{nazari2026rank} measure the effective rank of linear
attention hidden states and propose post-training rank pruning of
$K$ and $Q$ matrices. \citet{anonymous2026staterank} report
``state-rank stratification'' during pretraining: linear-attention
heads bifurcate into persistently low-rank and high-rank groups.
Both papers measure rank in the \emph{fast-weight memory} inside
an attention layer (a $d \times d$ accumulator), and both make
\emph{descriptive} claims about what trained networks end up with.
The object of study here is different --- the explicit per-position
matrix latents $\Z$ on the matrix-CODI feedback path, not a
fast-weight memory inside attention --- and the claim is a
mechanism claim about the training objective
(Proposition~\ref{prop:jac}) that we test by constructing four
nonlinear-in-$\Z$ positive controls.

\paragraph{Dynamics within latent CoT.}
\citet{anonymous2026dynamics} run multiple intervention protocols
on latent CoT hidden states (zero, mean, step-wise mean, Gaussian
noise) and an early-stop decoding that truncates latent
computation after step $k$. Their early-stop decoding is a cousin
of our rank-$k$ ablation on a different axis --- step depth vs.\
spectral truncation. Their step-wise causal structure is consistent
with our linear-probe decay $\Z[1]\to\Z[5]$ (early positions carry
information).

\paragraph{Rank decay from depth.}
\citet{dong2021attention} showed that pure attention loses rank
doubly exponentially with depth. Their subject is the rank of
activations across a stack of attention layers; ours is the rank
of an explicit matrix latent on a feedback path. The two phenomena
are distinct.

\paragraph{Implicit low-rank bias.}
Gradient descent on matrix-factorization losses has an implicit
bias toward low-rank solutions
\citep{gunasekar2017implicit,arora2019implicit,razin2020implicit},
and \citet{kobayashi2024weightdecay} show that weight decay
specifically induces low-rank attention products via an
equivalence with nuclear-norm regularization. These concern bias
through the parameter space; Proposition~\ref{prop:jac} concerns
the loss gradient through the readout. Our three-seed effective
rank spread $\{4,12,13\}$ under AdamW at weight decay $0.01$ is
consistent with an implicit low-rank attractor plus seed-dependent
convergence, and inconsistent with strong low-rank collapse (which
would predict all three seeds at the same low rank).

\paragraph{Alternative substrates and probe critique.}
\citet{wang2025latentvocab} argue that latent reasoning lives in
the vocabulary column space, not in the SVD directions of the
hidden state. If so, rank-$k$ truncation on $\Z$ is the wrong
observable; our negative result does not rule that out.
\citet{li2024optimal} show that zero/resample ablations (of which
rank-$k$ truncation is an instance) overestimate component
importance relative to optimal ablation; that cuts in our favor
(optimal ablation would make the curves flatter, not bumpier) but
is the right citation for the probe we use. The negative control
in \S\ref{sec:pc-negctrl} is the empirical analog: running
rank-$k$ on a model with no $\Z$ reproduces a flat curve, so the
probe alone cannot isolate rank-blindness from position-irrelevance.
